% ECONOMICS_PROBLEM: {"id": "D72-216", "title": "Breaking clientelistic equilibria", "jel_code": "D72", "source_id": "OP-0384"}
% !TeX program = xelatex
\documentclass[11pt,a4paper]{article}
\usepackage[margin=23mm]{geometry}
\usepackage{fontspec}
\setmainfont{Latin Modern Roman}
\usepackage{amsmath,amssymb,mathtools}
\usepackage{xcolor,enumitem,booktabs,longtable,array,xurl,hyperref}
\definecolor{muted}{HTML}{53636C}
\hypersetup{colorlinks=true,urlcolor=blue,linkcolor=blue}
\setlength{\parindent}{0pt}
\setlength{\parskip}{5pt}
\newcommand{\R}{\mathbb R}
\newcommand{\E}{\mathbb E}
\newcommand{\N}{\mathbb N}
\newcommand{\ind}{\mathbf 1}
\newcommand{\Var}{\operatorname{Var}}
\newcommand{\Cov}{\operatorname{Cov}}
\newcommand{\supp}{\operatorname{supp}}
\DeclareMathOperator*{\argmin}{arg\,min}
\DeclareMathOperator*{\argmax}{arg\,max}
\newcommand{\entrymeta}[3]{{\sffamily\small\color{muted}\textbf{\texttt{#1}}\quad|\quad #2\par #3\par}}
\newcommand{\Problemref}[1]{\texttt{#1}}
\newcommand{\JELref}[2]{\texttt{#2} (JEL \texttt{#1})}
\begin{document}
\section*{Breaking clientelistic equilibria}
\textbf{Problem ID:} \texttt{D72-216}\quad \textbf{JEL:} \texttt{D72}\par
% BEGIN ECONOMICS STATEMENT
\label{problem:OP-0384}
{\sffamily\small\color{muted}Original subject group: Political economy and institutions\par}
\label{definitions:problem:30007666}
\textbf{Audit classification:} Published research agenda. Three authors, title and 2026 NBER identifier checked. Interventions and future research are mentioned; the exact persistence test and strategic model are editorial, and chapter full text was not inspected.
\subsection*{Definitions and precise research target}
\textbf{Complete editorial problem.} Fix a political jurisdiction with politicians, brokers, and voters. Define clientelism as targeted material benefits exchanged for votes or electoral participation under specified monitoring arrangements. Let \(p\) be a feasible intervention changing benefit eligibility, ballot secrecy, broker monitoring, or public information. Let \(C_t(p)\) be the independently measured frequency or value of clientelistic exchange and \(G_t(p)\) public-good provision. Let \(V_t(p)\) measure responsiveness to programmatic performance rather than private transfers, using a predeclared empirical rule. Estimate \(\tau_C(T,p)=E[C_T(p)-C_T(0)]\), together with changes in \(G\) and \(V\), across multiple election cycles. Specify strategic responses by politicians and brokers, financing sources, and substitution into turnout buying, coercion, or patronage. Short-run reductions in reported vote buying do not establish a transition to a durable programmatic equilibrium. Identification requires credible intervention variation, independent behavior measurement, and explicit assumptions about spillovers and secrecy. Determine which bundles alter the incentives sustaining clientelistic exchange and whether effects survive withdrawal. The author's chapter abstract explicitly includes interventions and future research; the particular durability test and outcome vector here are an adaptation rather than an original published conjecture.

\textbf{Definitions and admissible scope.} A clientelistic exchange pairs targeted benefit \(b_i\) with an expected political act \(v_i\), under a specified monitoring or enforcement mechanism. Define whether turnout inducements, vote buying, coercion, and patronage are included or measured as separate substitutes. A political equilibrium is a strategy profile with sequentially specified information and no profitable unilateral deviation in the chosen solution concept; an observed frequency is not itself an equilibrium. A durable intervention effect means a stated reduction across \(K\ge2\) election cycles, including a cycle after withdrawal. Programmatic responsiveness can be a causal slope or a predeclared interaction between verified performance information and electoral support; define its assignment and measurement. Broker financing and strategic responses enter the policy regime. Fix election-cycle dates \(t=1,\ldots,K\), \(K\ge2\), a withdrawal cycle \(k_0<K\), and a positive baseline-jurisdiction weighting law. Define durability as \(E[C_t(p)-C_t(0)]\le-\epsilon\) for every declared postwithdrawal cycle, with \(\epsilon>0\) fixed. For programmatic responsiveness, classify incumbents' independently verified pretreatment performance as \(Q\in\{H,L\}\) under a fixed rule. Within each category randomize truthful disclosure \(M\in\{0,1\}\), where \(M=0\) withholds the message and \(M=1\) discloses that incumbent's actual category. With support outcome \(B_i(p,M)\), define \(V(p)=E[B_i(p,1)-B_i(p,0)\mid Q=H]-E[B_i(p,1)-B_i(p,0)\mid Q=L]\), using fixed eligible-voter weights and positive category probabilities. This is heterogeneity in disclosure effects; it does not randomize true performance, and cross-category causal comparisons require additional assumptions. A strategic model specifies players, feasible actions, monitoring, financing, timing, beliefs, and a selected equilibrium. Report the intervention identified set and whether the selected durability criterion holds.
\subsection*{Variants and assumption changes}
\begin{enumerate}
\item Monitoring and secrecy (source-motivated): vary ballot privacy or credible broker-monitoring limits and measure all forms of contingent political exchange, allowing substitutions into turnout buying or patronage.
\item Eligibility and fiscal access (editorial): replace discretionary benefits with rule-based eligibility, specifying funding and enforcement. Test whether reduced clientelism survives program withdrawal and whether a different broker network or coercion replaces it.
\end{enumerate}
\subsection*{References and source audit}
\textbf{Support and access.} Three authors, title and 2026 NBER identifier checked. Interventions and future research are mentioned; the exact persistence test and strategic model are editorial, and chapter full text was not inspected. \textbf{Locator:} Author research page, chapter abstract, last sentence.
\begin{enumerate}\raggedright
\item\hypertarget{definitions:source:30007666:1}{} Jimmy Graham; Horacio Larreguy; Pablo Querubín. 2026. \emph{Clientelism: How It Works, Why It Persists and How to Break It}. Research page, abstract under the exact chapter title, final sentence.
\url{https://sites.google.com/site/pabloquerubin/research}.
DOI: \url{https://doi.org/10.3386/w34761}.
\end{enumerate}

\subsection*{Common conventions: Source mathematical conventions}

These conventions apply to each entry unless explicitly replaced.

\textbf{Models and identification.} A primitive model \(m\) specifies all probability laws, feasible choices, information, equilibrium and measurement rules used by its target. The admissible class \(\mathcal M\) is fixed independently of outcomes. If \(P_O(m)\) is its observable probability law and \(\tau(m)\) the target, its identified set at observable law \(P\) is
\[
 \mathcal I_\tau(P)=\{\tau(m):m\in\mathcal M,\ P_O(m)=P\}.
\]
Point identification means that this set is a singleton. An empty set signals incompatibility of data and assumptions. Coordinate bounds need not describe the joint set. Bounds are sharp only if every retained target is attainable or arbitrarily approachable by compatible models. Finite bounds require explicit restrictions on unobserved quantities; unrestricted confounding, hidden wealth, extrapolation or arbitrary equilibrium selection can yield uninformative sets.

\textbf{Causal interventions.} A potential outcome \(Y_i(p)\) is the outcome under a complete policy regime \(p\), including timing, financing, information and market spillovers. The target population and units are fixed before treatment. With interference, use the complete assignment vector or a justified exposure mapping; individual no-interference notation alone cannot identify an economy-wide intervention. A difference of mean potential outcomes does not require the cross-world joint distribution. Conditioning on post-treatment migration, survival, enrollment or employment changes the estimand and needs separate justification.

\textbf{Statistical statements.} An estimator, confidence set, forecast or test specifies a sampling law, model class and loss/coverage criterion. Uniform \((1-\alpha)\)-coverage means \(\inf_{m\in\mathcal M}\Pr_m\{\tau(m)\in C_n\}\geq1-\alpha\), or an explicitly asymptotic version. The whole parameter space trivially has coverage, so informativeness requires a stated width, risk or power objective. A forecasting improvement is relative to a named benchmark, information set and evaluation population.

\textbf{Equilibrium and optimization.} An equilibrium comprises feasible plans, optimality, beliefs where needed, budget constraints and all market/resource clearing. The primitives or a stated benchmark define each correspondence \(\mathcal E(p;m)\). Nonemptiness is assumed or proved, never inferred just from compact policy sets. A selected-equilibrium target uses a declared selection rule; a robust target quantifies over all permitted equilibria. Use suprema or infima unless attainment is established. An \(\varepsilon\)-optimal policy is feasible and within \(\varepsilon\) of the finite optimum value. Report an empty feasible set separately.

\textbf{Welfare and accounting.} Normative utility, interpersonal normalization, social weights, numeraire, discounting and the use of public receipts are primitives. Behavioral utility need not coincide with the welfare criterion. Household welfare includes ownership income and rents once; adding profits again to compensating variation generally double-counts them. A monetary equivalent is defined by an explicit transfer and price convention, with monotonicity and range conditions. Average effects, mechanism contributions, transfers and welfare losses are different targets. Infinite sums require convergence and dynamic feasible plans require terminal or no-Ponzi conditions.

\textbf{Source versions.} A working-paper date, revision date and journal publication year are distinguished. A DOI for a journal version is not automatically the DOI of an earlier manuscript. Locators refer to the stated version and printed pages where specified. A metadata-only check does not verify a claim about a full-text conclusion. Every entry's source subsection narrows the provenance claim accordingly.
% END ECONOMICS STATEMENT
\end{document}
